\begin{theorem}
\label{thm:p01-real-pl-positive-current-candidate}
Let $K$ be a complete nontrivially valued non-Archimedean field, let $X$
be a good, separated, boundaryless, strictly $K$-analytic space, and let
$h:X\to\mathbf R$ be an $\mathbf R$-PL function.  The following conditions
are equivalent:
\begin{enumerate}
\item the Chambert--Loir--Ducros current
      $\dprimeddprime[h]$ is positive;
\item the residue class
      \[
      L_h(x)\in\operatorname{Pic}(\operatorname{red}(X,x))
                    \otimes_{\mathbf Z}\mathbf R
      \]
      is nef for every $x\in X$;
\item $h$ is semipositive in the sense of Gubler--Jell--Rabinoff.
\end{enumerate}
Consequently, if $X$ is allowed to have boundary, positivity of
$\dprimeddprime[h]$ on $\operatorname{Int}(X)$ is equivalent to
semipositivity of $h|_{\operatorname{Int}(X)}$.
\end{theorem}

\begin{proof}
The equivalence of (2) and (3) is
\cite[Proposition~7.1]{GJR26}, and (3) implies (1) by
\cite[Theorem~7.14]{GJR26}.  We prove that (1) implies (2).

Fix $x\in X$.  By \cite[Proposition~5.8 and Lemma~5.10]{GJR26}, after
shrinking to an analytic neighbourhood of $x$, write
\begin{equation}
\label{eq:p01-real-pl-decomposition}
                         h=\sum_{i=1}^{s}\lambda_i h_i,
\end{equation}
where the $h_i$ are integral PL functions and
$\lambda_i\in\mathbf R$.

We first construct one integral PL function that dominates the finitely
many $h_i$ in both current directions.  Choose a compact strictly analytic
neighbourhood $Y$ of $x$, with $x\in\operatorname{Int}(Y)$, on which all
$h_i$ are defined.  Since $Y$ is compact, the model-function theorem
\cite[\S5.6]{GJR26} represents each $h_i$ on a formal model of $Y$.
Passing to one formal model dominating this finite family, as in the proof
of \cite[Corollary~5.9]{GJR26}, gives a model $V_0$ of
$\operatorname{red}(X,x)$ carrying line bundles that represent all
$L_{h_i}(x)$.  Since $x$ is inner, $V_0$ is proper.  The model construction
in the proof of \cite[Proposition~7.10]{GJR26} now applies Chow's lemma to
produce a proper birational morphism $V\to V_0$ with $V$ projective.  Denote
the pullbacks of the representing line bundles by $M_i$.  Thus $V$ is one
projective model of $\operatorname{red}(X,x)$ on which all $L_{h_i}(x)$ are
represented.  Choose an ample line bundle $A$ on $V$.
The realization statement used in the same proof, namely
\cite[\S6.6.5]{CLD12}, produces an integral PL germ $a$ at $x$ such that
$(V,A)$ is a model of $L_a(x)$.  Shrinking once, regard $a$ and all the
$h_i$ as functions on the same neighbourhood.

For each $i$, the line bundles
\[
                         A^{\otimes N}\otimes M_i,
             \qquad     A^{\otimes N}\otimes M_i^{-1}
\]
are ample for all sufficiently large $N$.  Indeed, apply the ample-twist
criterion \cite[Tag~0892]{stacks-project} to the identity morphism of $V$,
the line bundles $M_i^{\pm1}$, and $A$.  Since the family is finite, fix one
positive integer $N$ that works for every $i$.  Thus
\begin{equation}
\label{eq:p01-ample-domination-residues}
                       NA+M_i\quad\text{and}\quad NA-M_i
\end{equation}
are nef on $V$.  In particular, $a$ and each $Na\pm h_i$ are semipositive
at $x$ by \cite[Proposition~7.1]{GJR26}.  The semipositive locus is open by
\cite[Remark~7.3(4)]{GJR26}; after intersecting the finitely many such loci,
all these functions are semipositive on one open neighbourhood $U$ of $x$.
Consequently, \cite[Theorem~7.14]{GJR26} gives positive currents on $U$
\begin{equation}
\label{eq:p01-current-domination}
 T_a:=\dprimeddprime[a]\geq0,
 \qquad NT_a+T_i\geq0,
 \qquad NT_a-T_i\geq0,
 \qquad T_i:=\dprimeddprime[h_i].
\end{equation}

Set $\varepsilon_m=1/m$.  Choose $q_i^{(m)}\in\mathbf Q$ so close to
$\lambda_i$ that, with $r_i^{(m)}=q_i^{(m)}-\lambda_i$,
\begin{equation}
\label{eq:p01-rational-error-budget}
                    N\sum_{i=1}^{s}|r_i^{(m)}|\leq\varepsilon_m.
\end{equation}
Define the $\mathbf Q$-PL function
\begin{equation}
\label{eq:p01-rational-positive-approximants}
                       g_m:=\sum_{i=1}^{s}q_i^{(m)}h_i
                                  +\varepsilon_m a.
\end{equation}
For a nonzero real number $r$, write $\sigma(r)=r/|r|$; when $r=0$, omit
the corresponding summand.  Linearity of residues and currents, the
assumed positivity of $T_h:=\dprimeddprime[h]$, and
\eqref{eq:p01-current-domination} give
\begin{align}
\label{eq:p01-positive-decomposition}
 \dprimeddprime[g_m]
  ={}&T_h
    +\left(\varepsilon_m-N\sum_i|r_i^{(m)}|\right)T_a \\
   &+\sum_{r_i^{(m)}\ne0}|r_i^{(m)}|
       \left(NT_a+\sigma(r_i^{(m)})T_i\right)\geq0
\end{align}
on $U$.  Notice that this argument tests positivity against every positive
compactly supported form on $U$; no choice of a tropical chart or skeleton
is involved.

Choose a positive integer $d_m$ such that $d_mg_m$ is integral PL.  Apply
\cite[Corollary~18.8.8]{CLD25} to the trivial line bundle with the integral
PL metric $\lVert1\rVert=\mathrm e^{-d_mg_m}$.  Its first Chern current is
$\dprimeddprime[d_mg_m]$.  We use the corollary in its printed scope: its
statement assumes a boundaryless analytic space and does not impose purity
or reducedness.  In view of
\eqref{eq:p01-positive-decomposition}, its residue at $x$ is nef.  Dividing
by $d_m$ shows that
\begin{equation}
\label{eq:p01-nef-approximating-residues}
             L_{g_m}(x)=\sum_iq_i^{(m)}L_{h_i}(x)
                              +\varepsilon_mL_a(x)
\end{equation}
is nef.

On the fixed projective model $V$, the class in
\eqref{eq:p01-nef-approximating-residues} is represented by
$\sum_iq_i^{(m)}M_i+\varepsilon_mA$.  By
\cite[Lemma~6.11]{GJR26}, this representative is nef.  Hence for every
proper integral curve $C\subset V$,
\[
  0\leq
  \sum_iq_i^{(m)}\deg_C(M_i)+\varepsilon_m\deg_C(A).
\]
Here $q_i^{(m)}\to\lambda_i$ by
\eqref{eq:p01-rational-error-budget}, and $\varepsilon_m\to0$.  Taking the
limit yields
\[
                       \sum_i\lambda_i\deg_C(M_i)\geq0.
\]
This holds for every proper integral curve on $V$, so
$\sum_i\lambda_iM_i$ is nef.  It represents
$L_h(x)=\sum_i\lambda_iL_{h_i}(x)$ by
\cite[Lemma~6.18]{GJR26}; hence $L_h(x)$ is nef.  Since $x$ was arbitrary,
(2) holds.

For an analytic space with boundary, its interior is a boundaryless open
analytic subspace.  Apply the statement already proved to
$\operatorname{Int}(X)$.  This proves the final assertion.
\end{proof}
